\begin{statement}{E.35}{Proved}{All retained joint fibres pass the order-three torsion-class test}
Let $T_{23}=T_{2,6}\times T_{3,6}$ be the exact joint
Ferrari target of E.24. Write
$e=a^2bc$, $z_4=(abc)^4$, and
$e'=a'^2b'c'$ for the image of the same word under a
retained braid triple. For every one of the $864$
points of E.24, and every compatible multiplier lift
$\lambda\equiv1\pmod{72}$, there is $h\in T_{23}$ with
\[
                         e'=he^\lambda h^{-1}.
 \tag{E.35.1}
\]
Thus every retained fibre passes E.34 already in the
actual joint target. On the ternary factor, all $216$
parameter points have been evaluated. There are $15$
distinct values of $e'$; for every such value,
\[
 e'z_4^j\sim_{T_{3,6}}e
              \quad\Longleftrightarrow\quad j\equiv0\pmod3,
             \qquad j\in\mathbb Z/9.
 \tag{E.35.2}
\]
All $45$ positive comparisons have saved literal
$B_4$ conjugators. The assertion concerns these exact
finite fibres and does not assert a full profinite lift.
\end{statement}
\begin{proof}
\par\medskip\noindent\textbf{Complete conjugacy in the ternary target.}\quad
Use the actual projection from E.23,
\[
 1\longrightarrow J\longrightarrow T_{3,6}
       \longrightarrow T_{3,2}\longrightarrow1,
       \qquad |T_{3,2}|=648.
 \tag{E.35.3}
\]
The complete four-strand Schreier image, retained in
E.23, gives a class-two exponent-three group $J$ with
\[
 \dim_{\mathbb F_3}\log J=30,\qquad
 \dim_{\mathbb F_3}\log[J,J]=18.
 \tag{E.35.4}
\]
The logarithm and exponential here are the exact
class-two maps on $J$ certified in E.23. We do not
apply a logarithm to the whole degree-six group.
On this kernel the BCH identity is
$u*v=u+v+\frac12[u,v]$, with $\frac12=2$ in
$\mathbb F_3$.

The new program independently enumerates all $648$
base elements using actual words in $a,b,c$, and
reconstructs the whole $30$-dimensional kernel from
the resulting Schreier relations and their commutators.
It verifies equality with the saved kernel of E.23
and records a literal $B_4$ word for every basis
element. Let $Z=\log[J,J]$, let
$V=\log J$, and put $A=\operatorname{Ad}(e)$ on $V$.
All these maps are computed by substitution in the
original finite group.

For a prescribed target $y$, enumerate every base
conjugator $q$ satisfying
$q\bar e q^{-1}=\bar y$ in $T_{3,2}$, and choose its
recorded lift $g_q$. There are $72$ base conjugates
of $\bar e$ and nine possible base conjugators for
each of them. Conjugacy above a given $q$ is equivalent
to conjugating $e$ to $g_q^{-1}yg_q$ by an element
of $J$. Its first equation is
\[
 (1-A)u\equiv
   \log\bigl(g_q^{-1}yg_qe^{-1}\bigr)\pmod Z.
 \tag{E.35.5}
\]
This is the exact equation modulo the central subgroup
$[J,J]$. Its homogeneous solution space is
\[
 V_1=\ker\bigl(V\xrightarrow{1-A}V/Z\bigr).
\]
The corresponding group $\exp(V_1)$ is the full
centralizer of $e$ modulo $[J,J]$. For $v\in V_1$,
$Av-v$ is central, so the exact class-two formula gives
\[
                    [\exp(v),e]=\exp((1-A)v).
 \tag{E.35.6}
\]
In particular the second equation, after any solution
$u$ of (E.35.5), is the linear equation on \emph{all}
of $V_1$ given by the remaining central discrepancy.
Multiplying $\exp(u)$ on the right by $\exp(v)$
corrects that discrepancy. Its full homogeneous fibre
is $\exp(V_2)$, where
\[
 V_2=\ker\bigl((1-A)|_{V_1}\bigr).
 \tag{E.35.7}
\]
The first solution set is an entire coset of
$\exp(V_1)$; the second is an entire coset of
$\exp(V_2)$. Thus both kernel spaces retain every
possible conjugator lift. An inconsistent equation
excludes that complete coset, and enumeration of all
base conjugators decides conjugacy in $T_{3,6}$.
The computed ranks and dimensions are
\[
 \begin{array}{c|ccc}
       &V&V_1&V_2\\\hline
 \dim&30&22&10
 \end{array},
 \qquad
 \operatorname{rank}\overline{(1-A)}=8,
 \quad\operatorname{rank}((1-A)|_{V_1})=12.
 \tag{E.35.8}
\]

Every ternary input is reconstructed from its retained
E.23 braid triple and its complete two-dimensional
homogeneous correction plane. For each of the $24$
indices, the program evaluates all nine parameters
and all nine central translates of the resulting
$e'$. The following exact facts are checked:
\[
 e^3=z_4,\quad |z_4|=9,\quad |e|=27,
 \qquad (e')^3=z_4.
 \tag{E.35.9}
\]
All nine parameters in a given indexed plane have
the same value of $e'$. Across the $24$ indices
there are $15$ values, giving $135$ distinct translated
targets. The complete two-stage algorithm gives
exactly the three positive translates $j=0,3,6$
for each value. The other six translates fail after
the complete base and kernel search. This proves
(E.35.2).

The resulting $45$ conjugators have literal braid
words of lengths between $6$ and $912$. They are
obtained by multiplying the Schreier words for a
logarithm's basis coordinates and then correcting
the resulting central BCH term with the recorded
central commutator words. Every reconstructed word
is evaluated in the actual group and checked against
the computed conjugator. A separate program imports
no conjugacy solver: it evaluates all $45$ words,
verifies the conjugacy identities, and reconstructs
all $216$ original parameter inputs independently.
The identity GT element, whose triple is $(a,b,c)$,
is retained as a literal control.

For the multiplier scope, the binary CRT corrections
of E.24 have identity image in the entire ternary
source, so this test covers both binary parameters
of each point. Both multiplier branches have the
same ternary triple; the exact braid-generator orders
divide nine. Moreover
\[
 e^{73}=e z_4^6,
 \qquad
 e^\lambda=e z_4^{(\lambda-1)/3}
          \quad(\lambda\equiv1\pmod{72}).
 \tag{E.35.10}
\]
The exponent $(\lambda-1)/3$ is a multiple of three.
All its residues modulo nine are among the three
positive translates in (E.35.2). Hence $e'$ is
conjugate to $e^\lambda$ in the ternary group itself
for every compatible higher multiplier lift.

An exhaustive auxiliary control tests the same
two-stage method on the class-two exponent-three
group with generators $a,b,c$ and central relation
$[b,c]=1$, of order $3^5=243$. The order-three
automorphism fixes $a$, sends $b$ to $b+c$ in
logarithmic coordinates, and fixes $c$; it sends
$[a,b]$ to $[a,b]+[a,c]$. All $243$ target values
are compared with the complete enumeration of all
$243$ conjugators. All decisions agree; the class
has nine elements and the two successive ranks are
one and one.

\par\medskip\noindent\textbf{The binary factor and the actual product.}\quad
Let $U=\ker(T_{2,6}\to\operatorname{SL}_2(\mathbb F_2))$.
It is the finite $2$-group of E.24, and its center
contains $\langle z_4\rangle$ of order four. In
$T_{2,6}/\langle z_4\rangle$, both $e^\lambda$ and
$e'$ have order three. Their images in
$\operatorname{SL}_2(\mathbb F_2)=S_3$ are the same
specified generator of $C_3$: the three retained
braid generators have the original degree-two images,
and, writing $A,B$ for their linear transvections,
\[
     \overline e=A^2BA=BA,
     \qquad\overline {e'}=BA,
     \qquad(BA)^\lambda=BA.
 \tag{E.35.11}
\]
Here $A^2=1$ and $\lambda\equiv1\pmod3$.

We recall an elementary proof of the needed
coprime conjugacy fact. Two order-three complements
to a normal finite $2$-group $P$ are conjugate by
$P$. Induct on $|P|$ and use the nonzero characteristic
elementary abelian subgroup $D=\Omega_1Z(P)$.
The induction hypothesis conjugates the two
complements to agree modulo $D$. For their generators
with the same specified image in $C_3$, write the
remaining difference as $t=vs$, $v\in D$, and let
$\alpha$ be conjugation by $s$. The equation $t^3=1$
is
\[
             v+\alpha v+\alpha^2v=0.
\]
In the $\mathbb F_2$-vector space $D$, take
$w=\alpha v$. Then
$(1-\alpha)w=v$, so conjugation by $w$ removes
the difference. This completes the induction and
keeps the generator orientation fixed.

Apply this argument to the preimage of the displayed
$C_3$ in $T_{2,6}/\langle z_4\rangle$, whose normal
$2$-group is $U/\langle z_4\rangle$. Lifting its
conjugator gives
\[
              e'=h_2e^\lambda h_2^{-1}z_4^j.
\]
Both cubes are $z_4^\lambda$: this follows from the
literal braid identity of E.34 and E.24's full-twist
calculation. Consequently $3j=0\pmod4$, so $j=0$.
This proves the binary equality in the group itself.

Finally E.24 proves that the actual joint target is
the full product $T_{2,6}\times T_{3,6}$. Pairing the
binary conjugator with the ternary one therefore
gives an element of that actual target and proves
(E.35.1), with no further joint-generation assumption.
The new executed files are
\texttt{de\_ferrari\_three\_conjugacy.py},
\texttt{de\_ferrari\_three\_torsion.py}, and
\texttt{de\_ferrari\_three\_torsion\_words.py},
with the complete centralizer matrices, translated
target decisions, indexed parameter map, and literal
conjugators.
\end{proof}
